\documentclass[UTF8,a4paper,11pt]{ctexart}
\usepackage[a4paper,margin=2.3cm]{geometry}
\usepackage{amsmath,amssymb,mathtools}
\usepackage{enumitem}
\usepackage[hidelinks]{hyperref}
\usepackage{fancyhdr}
\setlist[itemize]{leftmargin=2.2em,itemsep=0.25em,topsep=0.35em}
\setlength{\parindent}{2em}\setlength{\parskip}{0.35em}
\pagestyle{fancy}\fancyhf{}\fancyhead[L]{离散数学与结构（I）}\fancyhead[R]{Lec 10}\fancyfoot[C]{\thepage}\setlength{\headheight}{14pt}
\title{唯一分解整环、主理想整环与欧几里得整环}\author{zhangjunshuo}\date{2026 年 10 月}
\begin{document}
\maketitle
\begin{quote}
本讲根据 2026 年 10 月 15 日的三段课堂转写、手写原稿和课程教材整理。课堂的主线是把整数"唯一分解"推广到一般整环：先厘清单位、整除与相伴，再定义素元与不可约元，然后说明唯一分解整环等价于"因子链条件 $+$ 素性条件"，并沿欧几里得整环 $\to$ 主理想整环 $\to$ 唯一分解整环这条链下来说明常见环都满足这些性质；最后转向用理想构造域，并引出分裂域。

每个分节开头都先交代本节用到的记号；所有符号在首次使用处给出定义，不依赖读者记得前面讲次的写法。
\end{quote}
\tableofcontents\clearpage

\section{单位、整除与相伴}


本讲从整数"唯一分解"这件事出发，把它推广到一般整环。推广时有两件事必须先说清楚：\textbf{哪些元素算"可逆"}（单位），以及\textbf{两个元素在什么意义下算"同一个因子"}（相伴）。这一节把这两个概念和整除的关系理清。

\textbf{本节用到的记号，先一次交代清楚：}

\begin{itemize}
\item $R$：整环（交换含幺环且 $1\ne0$、无零因子）；$R^\times$：$R$ 的单位群，即全体可逆元。
\item $a\mid b$：$a$ 整除 $b$，即存在 $c\in R$ 使 $b=ac$。
\item $a\nmid b$：$a$ 不整除 $b$。
\item $(a)=\{ra:r\in R\}$：由 $a$ 生成的主理想。
\item $\mathbb Z[\sqrt{d}]=\{a+b\sqrt{d}:a,b\in\mathbb Z\}$（$d$ 为无平方因子的整数）：$\mathbb Z$ 添上 $\sqrt{d}$ 得到的环，是本讲反复使用的反例来源。
\end{itemize}


\subsection{单位与整除}


含幺环 $R$ 的\textbf{单位群}是

\[
R^\times=\{u\in R:\exists v\in R,\ uv=vu=1\},
\]

即全体有乘法逆元的元素。它在乘法下构成群。$\mathbb Z$ 中 $\mathbb Z^\times=\{\pm1\}$；域 $F$ 中 $F^\times=F\setminus\{0\}$；而 $F[x]$ 中 $(F[x])^\times=F\setminus\{0\}$（非零常数），见 [Lec 09 §4](../lec09/sections/04-units-fractions-characteristic.md)。

\textbf{整除}的定义是

\[
a\mid b\iff \exists c\in R,\ b=ac .
\]

整除有两条常用性质，都由定义直接得到：

\begin{itemize}
\item 传递性：$a\mid b$ 且 $b\mid c$ $\Longrightarrow$ $a\mid c$；
\item 与理想的关系：$a\mid b\iff (b)\subseteq(a)$。\textbf{"整除"是"包含"的反向写法}，这一点在下一节讨论因子链时非常关键。
\end{itemize}


\subsection{相伴：什么时候两个元素算"同一个因子"}


若存在单位 $u\in R^\times$ 使 $a=ub$，就说 $a$ 与 $b$ \textbf{相伴}，记作 $a\approx b$。例如在 $\mathbb Z$ 中 $2\approx-2$，在 $F[x]$ 中 $f\approx cf$（$c\ne0$）。

相伴是一个等价关系（单位群是群，验证自反、对称、传递都用群公理），所以"相伴"把环中的元素分成若干类。为什么要引入它？因为在整环里"$a$ 和 $b$ 互相整除"恰恰等价于"$a$ 和 $b$ 相伴"，而\textbf{不是}等价于 $a=b$：

\textbf{命题。} 设 $R$ 是整环，$a,b\ne0$。则

\[
a\mid b\ \text{且}\ b\mid a\iff a\approx b .
\]

\textbf{证明。} （$\Leftarrow$）若 $a=ub$（$u$ 为单位），则 $b=u^{-1}a$，于是 $a\mid b$ 且 $b\mid a$。

（$\Rightarrow$）设 $b=ac$、$a=bd$（$c,d\in R$）。代入得 $a=acd$，即 $a(1-cd)=0$。因为 $a\ne0$ 且 $R$ 是整环（无零因子），所以 $cd=1$，即 $c,d$ 都是单位。于是 $b=ac$ 且 $c\in R^\times$，故 $a\approx b$。$\square$

\textbf{这条命题为什么重要。} 在 $\mathbb Z$ 里 $6=2\cdot3=(-2)\cdot(-3)$，如果"唯一分解"要求分解式逐字相同，那 6 就有两种分解。但 $2\approx-2$、$3\approx-3$，两式其实只差单位。所以唯一分解的正确说法必须是\textbf{"在相伴意义下唯一"}，这正是下一节要定义的\textbf{唯一分解整环}（Unique Factorization Domain，简称 \textbf{UFD}）里要写进去的。

\subsection{相伴的三种等价刻画}


在整环中以下三件事等价，可以互换使用：

\[
a\approx b\iff a\mid b\ \text{且}\ b\mid a\iff (a)=(b).
\]

第一、二条的等价就是上面的命题；第一条与第三条的等价来自"$a\mid b\iff(b)\subseteq(a)$"：互相整除就是两个主理想互相包含，也就是相等。

\textbf{读题提示：} 看到"$a$ 与 $b$ 相伴"，可以立刻换成"$(a)=(b)$"来算，往往更方便；反过来看到两个主理想相等，也可以换成相伴来谈元素。

\subsection{参考对照}


\begin{itemize}
\item [DF] §8.1（整除、相伴、不可约元与素元的一般讨论）。
\item [DS] §10.1（整环与单位；"相伴"只在 §8.1.2 命题 (vi) 里对 $\mathbb Z$ 定义过）。
\item 课程讲义 §7.3（相伴的定义 7.19）。
\item 原始材料：本地归档（手写原稿与录音转写），不纳入课程仓库。
\end{itemize}

\section{素元、不可约元与唯一分解整环}


上一节说明了"同一个因子"要在相伴意义下理解。这一节定义\textbf{不能再分解的元素}（不可约元）与\textbf{具有素数性质的元素}（素元），并用它们写出唯一分解整环的定义；随后给出两个不是唯一分解整环的经典反例。

\textbf{本节用到的记号：} $p,q$ 表示待考察的元素；$a\mid b$ 为整除；$a\approx b$ 为相伴（上一节）；$N$ 为范数，在 2.4 小节定义；$\mathbb Z[\sqrt{d}]$ 的定义见上一节。

\subsection{素元与不可约元}


设 $R$ 是整环，$p\in R$ 非零且\textbf{不是单位}。称 $p$ 是

\begin{itemize}
\item \textbf{素元}（prime），若
\end{itemize}

\[
  p\mid ab\ \Longrightarrow\ p\mid a\ \text{或}\ p\mid b\qquad(\forall a,b\in R);
\]
\begin{itemize}
\item \textbf{不可约元}（irreducible），若 $p=ab$ 只能有平凡分解：
\end{itemize}

\[
  p=ab\ \Longrightarrow\ a\ \text{是单位或}\ b\ \text{是单位}.
\]

两者的直觉不同：\textbf{素元说的是"整除一个乘积就能整除其中一个因子"（和素数一样），不可约元说的是"自己不能再拆成两个非单位的乘积"。} 在 $\mathbb Z$ 中两者恰好一致（$\pm2,\pm3,\pm5,\dots$），所以我们平时不会注意它们的区别；但在一般的整环里它们\textbf{不}等价。

\subsection{素元一定不可约}


\textbf{命题。} 在整环中，素元必为不可约元。

\textbf{证明。} 设 $p$ 素，且 $p=ab$。因为 $p\mid ab=p$，由素性 $p\mid a$ 或 $p\mid b$。不妨设 $p\mid a$，即 $a=pc$。代入 $p=ab$ 得 $p=pcb$，即 $p(1-cb)=0$。$p\ne0$ 且无零因子，故 $cb=1$，即 $b$ 是单位。所以 $p=ab$ 的分解必是平凡的，$p$ 不可约。$\square$

\textbf{反方向不成立}，这正是"唯一分解"这条线最要紧的坑。下一小节的反例给出一个不可约但不素的元素。反过来，"不可约 $\Rightarrow$ 素"在什么环里成立，是这条线的核心问题：

\begin{itemize}
\item 在\textbf{任意整环}里只有"素 $\Rightarrow$ 不可约"这一边（[DF] §8.3 Proposition 10）；
\item "不可约 $\Rightarrow$ 素"在 \textbf{PID} 里成立（[DF] §8.3 Proposition 11），在 \textbf{UFD} 里也成立（[DF] §8.3 Proposition 12）——后者其实与 UFD 的定义等价，见第 3 节的素性条件。
\end{itemize}


\textbf{注意不要循环论证：} 不能用"PID 是 UFD"去推"PID 里不可约 $\Rightarrow$ 素"，因为教材正是先用后者去证明前者的（[DF] §8.3 p.286 明确说明 Proposition 11 要排在 Theorem 14 之前使用）。

\subsection{唯一分解整环的定义}


称整环 $R$ 是\textbf{唯一分解整环}（Unique Factorization Domain，简称 UFD），若满足：

1. \textbf{存在性}：每个非零非单位元 $a\in R$ 都能写成有限多个不可约元的乘积
\[
   a=u\,p_1p_2\cdots p_n,\qquad u\in R^\times,\ p_i\ \text{不可约};
\]
2. \textbf{唯一性}：这样的分解在\textbf{相伴与次序}意义下唯一，即若
\[
   a=u\,p_1\cdots p_n=v\,q_1\cdots q_m
\]
   是另一个不可约分解，则\textbf{先有因子个数相等 $n=m$}，然后才谈重排：存在下标的重排使 $p_i\approx q_{\sigma(i)}$（$i=1,\dots,n$）。

"$n=m$"这一步不能省——它说的是两个分解的\textbf{长度}必须先一致，唯一性才可能成立。

\textbf{用一句话概括：UFD 就是"每个非零非单位元都有唯一的不可约分解"的整环。} 手写稿上写的正是这个意思——"every $r\ne0$ has unique factorization $r=u\,a_1\cdots a_t$，其中 $u$ 是单位、$a_i$ 不可约"。

定义里为什么要写一个单位 $u$ 在前面？因为 $a$ 本身可能就是一个不可约元（此时 $n=1$），或者 $a$ 是单位（此时 $n=0$，只剩下 $u$）；把单位单独提出去，两种情况都能统一写成上面的形状。\textbf{"在每个 $p_i$ 前也可以塞单位"这一点不影响唯一性，因为塞进去的单位可以全部并到最前面的 $u$ 里。}

对称地，$\mathbb Z$ 是 UFD，这正是算术基本定理；$F[x]$ 也是 UFD（$F$ 为域），这是多项式的唯一分解。

\subsection{反例一：$\mathbb Z[\sqrt{-5}]$ 不是 UFD}


取

\[
\mathbb Z[\sqrt{-5}]=\{a+b\sqrt{-5}:a,b\in\mathbb Z\},
\]

它在复平面上就是全体整点 $(a,b)$（横轴为实部、纵轴为 $\sqrt5$ 的倍数）。定义\textbf{范数}

\[
N(a+b\sqrt{-5})=a^2+5b^2\in\mathbb Z_{\ge0},
\]

它满足可乘性 $N(xy)=N(x)N(y)$（直接展开即可验证）。注意 $N(x)=0$ 只对 $x=0$ 成立，且

\[
x\ \text{是单位}\iff N(x)=1\iff x=\pm1 .
\]

原因是 $N(x)=a^2+5b^2=1$ 只有 $(a,b)=(\pm1,0)$。

\textbf{关键观察：$\mathbb Z[\sqrt{-5}]$ 中非零元素的范数最小值是 $1$。} 范数把元素映到非负整数，所以范数严格小于 $N(x)$ 的元素只有有限多个，判断 $x$ 是否可约时逐个排除即可。这就是课堂说的"画一个圆，圆内没有因子就不可约"。

\textbf{这一步的逻辑要说准。} 仅有"最小非零范数为 $1$"\textbf{本身不足以}断定 $x$ 不可约；必须再加上"$N(x)$ 的真因子都不落在范数的像里"。对 $\mathbb Z[\sqrt{-5}]$，范数的像是

\[
\{a^2+5b^2\}=\{0,1,4,5,6,9,10,14,16,21,\dots\},
\]

$2$ 与 $3$ 不在其中，所以 $N(x)\in\{4,9,6\}$ 的那些元素不可能有范数 $2$ 或 $3$ 的真因子，才推出不可约。写题时这两步都要写。

（另外说明出处：这个"画圆/取最近格点"式的论证在教材里是用来证 \textbf{$\mathbb Z[i]$ 是欧几里得整环}的（[DF] §8.1），教材对 $\mathbb Z[\sqrt{-5}]$ 的不可约性走的是范数路线、并把它列为习题；这里按课堂的讲法整理，但补足了上面这一步。）

于是可以验证：

\begin{itemize}
\item $N(2)=4$。若 $2=xy$ 是非平凡分解，则 $N(x),N(y)>1$ 且 $N(x)N(y)=4$，只能是 $N(x)=N(y)=2$。但 $a^2+5b^2=2$ 无整数解，所以 \textbf{$2$ 不可约}。
\item $N(3)=9$。同理需要 $N(x)\in\{3\}$，而 $a^2+5b^2=3$ 无整数解，所以 \textbf{$3$ 不可约}。
\item $N(1+\sqrt{-5})=1+5=6$。若可约则需 $N(x)\in\{2,3\}$（或 $6$ 的真因子 $1,2,3,6$ 中大于 $1$ 小于 $6$ 的），同样无解，所以 \textbf{$1\pm\sqrt{-5}$ 都不可约}。
\end{itemize}


但

\[
6=2\cdot3=(1+\sqrt{-5})(1-\sqrt{-5}),
\]

两个分解都用了不可约元，而且 $2,3$ 与 $1\pm\sqrt{-5}$ 互不相伴（范数分别是 $4,9$ 与 $6$，相伴的元素范数相同）。\textbf{所以 $\mathbb Z[\sqrt{-5}]$ 不是 UFD。}

\textbf{同一个环还有另一个同样有效的反例：} $9=3\cdot3=(2+\sqrt{-5})(2-\sqrt{-5})$，其中 $3$ 与 $2\pm\sqrt{-5}$ 都不可约（范数都是 $9$）且互不相伴。两个反例都对，只是来源不同——[DF] §8.3 用 $6$，课程讲义用 $9$。写作业时挑一个即可，\textbf{不要把两个混成一个式子}。

顺带得到"不可约 $\Rightarrow$ 素"的反例：在上式中 $2\mid(1+\sqrt{-5})(1-\sqrt{-5})$，但

\[
2\nmid 1+\sqrt{-5},\qquad 2\nmid 1-\sqrt{-5}
\]

（$a^2+5b^2$ 为偶数时 $a,b$ 一奇一偶，此时 $a,b$ 不可能同时为偶数，故 $2$ 除不尽）。\textbf{$2$ 不可约却不素。}

\subsection{反例二：$\mathbb Z[\sqrt{10}]$，以及为什么范数要换个方式用}


同样的想法对 $\mathbb Z[\sqrt{10}]=\{a+b\sqrt{10}:a,b\in\mathbb Z\}$ 也成立，但这里范数

\[
N(a+b\sqrt{10})=a^2-10b^2
\]

\textbf{可以取负值}，而且这个环的\textbf{单位群是无限的}：

\[
N(3+\sqrt{10})=9-10=-1,
\]

所以 $3+\sqrt{10}$ 是单位，它的幂给出无穷多个单位（它们对应 Pell 方程 $a^2-10b^2=\pm1$ 的无穷多组解）。因此在 $\mathbb Z[\sqrt{10}]$ 里\textbf{没有"离原点最近的正距离"}这回事——把元素乘上合适的单位可以让它任意接近 $0$，课堂说的"这些元素在实轴上稠密"就是这个意思。

于是改用\textbf{范数的绝对值}来数：因为 $|N(xy)|=|N(x)|\cdot|N(y)|$，判断 $x$ 不可约仍然只要看 $|N|$ 的真因子。逐个检查：

\begin{itemize}
\item $|N(2)|=4$。真因子只有 $2$，而 $a^2-10b^2=\pm2$ 无整数解，所以 \textbf{$2$ 不可约}。
\item $|N(3)|=9$。真因子 $3$，而 $a^2-10b^2=\pm3$ 无整数解，所以 \textbf{$3$ 不可约}。
\item $N(2+\sqrt{10})=4-10=-6$，$|N|=6$。真因子 $2,3$，都无解，所以 \textbf{$2\pm\sqrt{10}$ 不可约}。
\end{itemize}


而

\[
6=2\cdot3=(2+\sqrt{10})(-2+\sqrt{10}),
\]

同样是两种不同的不可约分解，所以 \textbf{$\mathbb Z[\sqrt{10}]$ 也不是 UFD}。而且 $2\mid 6=(2+\sqrt{10})(-2+\sqrt{10})$，但 $2\nmid 2\pm\sqrt{10}$（否则 $2$ 整除虚部 $\pm1$，不可能），所以 $2$ 同样是\textbf{不可约但不素}。

\textbf{小结这两个例子的区别：} $\mathbb Z[\sqrt{-5}]$ 的范数恒非负、最小值有正的下界（为 $1$），所以"画圆排除因子"直接可用；$\mathbb Z[\sqrt{10}]$ 的单位群无限，范数会变号，必须换成 $|N|$ 并用整除性排除。这也是为什么教材处理虚二次域比实二次域省事。

\subsection{参考对照}


\begin{itemize}
\item [DF] §8.1（不可约元与素元）、§8.3（UFD 的定义与例子，含 $\mathbb Z[\sqrt{-5}]$）。
\item [DS] §10.1（整环）。该书的"不可约 vs 素"只有 §8.3.2 对 $\mathbb Z$ 的一句说明，没有一般环的定义。
\item 课程讲义 §7.3（不可约与素元的定义 7.19、UFD 的定义 7.22）。
\item $\mathbb Z[\sqrt{10}]$ 这个反例三本教材都没有，本节按课堂讲法自行整理，范数与不可约性都逐条验证过。
\item 原始材料：$../../../课堂手记和录音转写/Lec10/$。
\end{itemize}

\section{因子链条件、主理想整环与欧几里得整环}


上一节的两个反例说明"存在不可约分解"并不自动保证唯一。这一节把 UFD 拆成\textbf{两个可以分别验证的条件}，然后说明主理想整环与欧几里得整环都自动满足它们。

\textbf{本节用到的记号：} $(a)$ 为主理想；$\subseteq$ 为集合包含；$N$ 为范数；$\deg f$ 为多项式次数。

\subsection{UFD 的两个等价条件}


\textbf{定理（UFD 的两个等价条件）。} 整环 $R$ 是 UFD，当且仅当它同时满足下面两条：

（\textbf{出处提示：} 这条等价刻画在 [DF] 里\textbf{没有}以定理形式出现——DF 把"存在性"与"唯一性"分别塞进 §8.3 Theorem 14（PID $\Rightarrow$ UFD）的证明里。给出这条等价定理的是课程讲义 §7.3 的定理 7.5 与定理 7.6（两个方向各一条）。所以本节这一条以讲义为准，[DF] 只提供两个条件的用法。）

\textbf{(F) 因子链条件（升链条件，ACC）。} $R$ 中不存在无穷序列

\[
a_1,\ a_2,\ a_3,\ \dots
\]

使每个 $a_{i+1}$ 都是 $a_i$ 的\textbf{真因子}（即 $a_{i+1}\mid a_i$ 且 $a_{i+1}\not\approx a_i$）。

\textbf{(P) 素性条件。} $R$ 中"不可约"与"素"等价，即每个不可约元都是素元。

（讲义把素性条件写成\textbf{双向}，其中"素 $\Rightarrow$ 不可约"在任意整环都成立（第 2 节已证），所以真正有内容的是"不可约 $\Rightarrow$ 素"这个方向。两个方向合起来才是"等价"。）

\textbf{为什么恰好是这两条？} 分别对应 UFD 定义里的"存在性"和"唯一性"：

\begin{itemize}
\item (F) 保证存在性。若存在无穷真因子链 $a_1,a_2,\dots$，则 $a_1$ 永远拆不完，得不到有限多个不可约元的乘积。反过来，有 (F) 就能反复取真因子，直到取不动为止——取不动的就是不可约元，于是每个元素都有不可约分解。
\item (P) 保证唯一性。唯一性的证明与整数情形逐字相同：设有两个不可约分解，取左边的 $p_1$。因为 $p_1$ 不可约，由 (P) 它是素元；又 $p_1$ 整除右边的乘积 $q_1\cdots q_m$，于是 $p_1$ 整除某个 $q_j$。重排下标设 $j=1$，即 $q_1=p_1c$；由于 $q_1$ 也不可约而 $p_1$ 不是单位，$c$ 必是单位，所以 $p_1\approx q_1$。两边约去（整环中可约分）后对剩余部分递归，就得到两个分解在相伴与次序下相同。
\end{itemize}


\textbf{用因子链条件的写法复述证明。} 上一段"$p_1$ 整除某个 $q_j$"这一步用的正是素性；如果没有 (P)，两个分解就可能像 $6=2\cdot3=(1+\sqrt{-5})(1-\sqrt{-5})$ 那样完全对不上——那里 $2$ 不可约但不素，所以它除得尽乘积却除不尽任何一个因子。

\subsection{主理想整环（PID）}


称\textbf{整环} $R$ 是\textbf{主理想整环}（Principal Ideal Domain，简称 PID），若它的每个理想都是主理想，即

\[
\forall I\trianglelefteq R,\ \exists a\in R,\ I=(a).
\]

\textbf{术语陷阱。} 有的讲义只要求"每个理想都是主理想"，不要求环是整环，并把这样的环也叫 PID。这时\textbf{"主理想环 $\Rightarrow$ UFD"是错的}：$\mathbb Z/6\mathbb Z$ 的每个理想都是主理想，但它不是整环，根本谈不上唯一分解（连零因子都有）。所以引用主理想环的结论时\textbf{必须补上"整环"前提}；本讲的 PID 一律指主理想\textbf{整环}。

\textbf{定理。} PID $\Longrightarrow$ UFD。

\textbf{证明。} 按上面的定理，只需分别验证 (F) 和 (P)。

\textbf{(F) 的验证。} 反设存在真因子链

\[
a_1,\ a_2,\ a_3,\ \dots,\qquad a_{i+1}\mid a_i,\quad a_{i+1}\not\approx a_i .
\]

由"$a_{i+1}\mid a_i\iff(a_i)\subseteq(a_{i+1})$"，对应得到理想的升链

\[
(a_1)\subsetneq(a_2)\subsetneq(a_3)\subsetneq\cdots
\]

（严格包含：因为 $a_{i+1}\not\approx a_i$，两主理想不相等）。令

\[
I=\bigcup_{i\ge1}(a_i).
\]

\textbf{$I$ 是理想}：对加法封闭，因为任意两个元素都落在某个 $(a_N)$ 里（取下标较大者）；对乘法封闭，因为 $r\in R$、$x\in(a_i)$ 时 $rx\in(a_i)\subseteq I$。

因为 $R$ 是 PID，$I=(b)$ 对某个 $b\in R$ 成立。而 $b\in I=\bigcup_i(a_i)$，所以 $b\in(a_N)$ 对某个 $N$ 成立，即 $(b)\subseteq(a_N)$。另一方面 $(a_N)\subseteq I=(b)$。于是 $(a_N)=(b)=I$，进而对一切 $i\ge N$ 有 $(a_i)\subseteq I=(a_N)\subseteq(a_i)$，即 $(a_i)=(a_N)$。这与链的严格性矛盾。$\square$

\textbf{(P) 的验证。} 设 $p$ 是 PID 中的不可约元，$p\mid ab$。考虑由 $p$ 与 $a$ 生成的理想 $(p,a)=(p)+(a)$。因为 $R$ 是 PID，存在 $s$ 使

\[
(p,a)=(s).
\]

于是 $p\in(s)$，即 $p=sc$ 对某个 $c$。因为 $p$ 不可约，$s$ 与 $c$ 中至少有一个是单位：

\begin{itemize}
\item \textbf{$s$ 是单位}：则 $(s)=R$，特别地 $1\in(p,a)$，于是存在 $\alpha,\beta\in R$ 使
\end{itemize}

\[
  \alpha p+\beta a=1 .
\]
  两边乘 $b$：$\alpha pb+\beta ab=b$。左边两项都被 $p$ 整除（第一项显然，第二项因为 $p\mid ab$），所以 $p\mid b$。
\begin{itemize}
\item \textbf{$c$ 是单位}：则 $s\approx p$，于是 $(s)=(p)$，即 $(p,a)=(p)$，从而 $a\in(p)$，即 $p\mid a$。
\end{itemize}


两种情形都得到 $p\mid a$ 或 $p\mid b$，所以 $p$ 是素元。$\square$

\subsubsection{推论：PID 中"不可约"与"极大理想"是一回事}


上面 (P) 的验证只用到 $p$ 不可约，把那个论证单独提出来就是一条很有用的判据：

\textbf{命题。} 设 $R$ 是 PID，$p\in R$ 非零非单位。则

\[
p\ \text{不可约}\iff(p)\ \text{是极大理想}\iff R/(p)\ \text{是域}.
\]

\textbf{证明。} （$\Rightarrow$）设 $p$ 不可约，且 $(p)\subseteq(s)\subseteq R$。因为 $R$ 是 PID，中间的理想由某个 $s$ 生成，且 $p\in(s)$ 即 $p=sc$。由 $p$ 不可约，$s$ 与 $c$ 中至少一个是单位：$s$ 是单位给出 $(s)=R$，$c$ 是单位给出 $s\approx p$ 从而 $(s)=(p)$。所以 $(p)$ 与 $R$ 之间没有别的理想。

（$\Leftarrow$）设 $(p)$ 极大，且 $p=ab$ 而 $a$ 不是单位。则 $(p)\subsetneq(a)$（严格：$p=ab$ 说明 $p\in(a)$，而 $a$ 不是单位说明 $a\notin(p)$），与极大性矛盾。所以这样的分解不存在，$p$ 不可约。

最后一步"极大 $\iff$ $R/(p)$ 是域"是上一节/下一节的定理（$R$ 交换含幺时极大理想给出域）。$\square$

\textbf{这条命题的用处：} 想造域时不必直接验证"没有中间理想"，只要找一个不可约元 $p$，$R/(p)$ 自动是域。

\subsection{欧几里得整环（ED）}


称整环 $R$ 是\textbf{欧几里得整环}（Euclidean Domain，简称 ED），若存在函数（\textbf{范数}）

\[
N:R\setminus\{0\}\longrightarrow\mathbb Z_{\ge0}
\]

使得对任意 $a,b\in R$、$b\ne0$，都存在 $q,r\in R$ 满足

\[
a=qb+r,\qquad r=0\ \text{ 或 }\ N(r)<N(b).
\]

\textbf{范数的约定要小心，各教材不同：}

\begin{itemize}
\item [DF] §8.1 的定义最宽：只要求 $N:R\to\mathbb Z_{\ge0}$ 且 $N(0)=0$。它\textbf{不要求}非零元的范数为正，甚至允许 $N\equiv0$——此时域自动成为欧几里得整环。
\item 课程讲义定义 7.21 写的是双向条件 $N(x)=0\iff x=0$，\textbf{比 [DF] 强}。后果是：若取 $N(f)=\deg f$，非零常数多项式的范数为 $0$，于是按讲义的口径 $F[x]$ \textbf{反而不是}欧几里得整区。这是讲义的一处实质缺陷，本节按 [DF] 的口径写（并把 $F[x]$ 的范数取成 $\deg f+1$ 以同时满足两种口径）。
\item 还有教材允许 $N$ 取到 $\mathbb Z$ 或额外要求 $N(x)\le N(xy)$，这些差别不影响主要结论。
\end{itemize}


这条性质就叫\textbf{带余除法}：$q$ 是商、$r$ 是余数，余数在范数意义下比除数"小"。注意条件的量词顺序：\textbf{先固定除数 $b\ne0$，再对任意被除数 $a$ 要求存在 $q,r$}。

\textbf{例子。}

\begin{itemize}
\item $R=\mathbb Z$，$N(a)=|a|$：这就是整数的带余除法。
\item $R=F[x]$（$F$ 为域），取 $N(f)=\deg f+1$（加 $1$ 是为了让零多项式的"范数"为 $0$，与上面 $N(0)=0$ 的约定一致）：这就是多项式的带余除法。任取 $f\ne0$ 与 $g$，做多项式除法得 $g=qf+r$，其中 $r=0$ 或 $\deg r<\deg f$，即 $N(r)<N(f)$。\textbf{注意这里要求 $F$ 是域}：只有除式首项系数可逆时才能做除法，例如 $\mathbb Z[x]$ 中就无法用 $2x$ 去除 $x$。
\item 任何\textbf{域} $F$：取 $N\equiv1$，则任取 $b\ne0$ 都有 $a=(ab^{-1})b+0$，余数为 $0$。所以域都是 ED。
\end{itemize}


\textbf{定理。} ED $\Longrightarrow$ PID。

\textbf{证明。} 设 $I\trianglelefteq R$。若 $I=(0)$ 已是主理想。否则 $I$ 含有非零元，在 $I\setminus\{0\}$ 中取一个范数 $N$ \textbf{最小}的元素 $a$（$N$ 取值于 $\mathbb Z_{\ge0}$，非空集合必有最小元）。

断言 $I=(a)$。包含 $\subseteq$ 显然。反之取 $b\in I$，用带余除法除以 $a$：

\[
b=qa+r,\qquad r=0\ \text{ 或 }\ N(r)<N(a).
\]

因为 $b\in I$ 且 $qa\in I$，所以 $r=b-qa\in I$。若 $r\ne0$，则 $r$ 是 $I$ 中范数比 $a$ 更小的非零元，与 $a$ 的取法矛盾。故 $r=0$，即 $b=qa\in(a)$。于是 $I=(a)$。$\square$

\subsection{一条链，以及它为什么不是等号链}


把上面的结论连起来，手写稿上写的正是这条链（箭头方向是"左边是右边的特例"）：

\[
\text{交换环}\ \supsetneq\ \text{整环}\ \supsetneq\ \text{UFD}\ \supsetneq\ \text{PID}\ \supsetneq\ \text{ED}\ \supsetneq\ \text{域}.
\]

读这条链要\textbf{特别小心方向}：它用的是反包含符号 $\supsetneq$，意思是越往右条件越强、越特殊。手写稿上把箭头写成"$\to$"，如果读成"蕴含"就对了（域 $\Rightarrow$ ED $\Rightarrow$ PID $\Rightarrow$ UFD），但\textbf{不能读成"ED 一定是域"}——那是错的，$\mathbb Z$ 就是欧几里得整环而不是域。正确的理解是：

\[
\text{域}\Longrightarrow\text{ED}\Longrightarrow\text{PID}\Longrightarrow\text{UFD}\Longrightarrow\text{整环}\Longrightarrow\text{交换环}.
\]

一个 ED 什么时候真的是域？当且仅当它的每个非零元都可逆，即它本身就是除环；$\mathbb Z$ 显然不满足（$2$ 不可逆），所以 $\mathbb Z$ 是"严格的" ED。

课堂上说"一个比一个更特殊、性质更好"，指的就是从右往左每走一步条件更强、环更特殊。这些包含关系都是\textbf{严格}的，反例可以记住：


\begin{center}
\begin{tabular}{lll}
关系 & 反例 & 为什么\\ \hline
整环 $\supsetneq$ UFD & $\mathbb Z[\sqrt{-5}]$ & 上一节的 $6$ 有两种不可约分解\\
UFD $\supsetneq$ PID & $\mathbb Z[x]$ & 它是 UFD，但 $(2,x)$ 不是主理想\\
PID $\supsetneq$ ED & $\mathbb Z\!\left[\frac{1+\sqrt{-19}}{2}\right]$ & 是 PID 但不是 ED（经典例子，本课不证）
\end{tabular}
\end{center}


\textbf{判断流程小结：} 遇到一个具体的环，先看能不能做带余除法（ED），能则它同时是 PID 和 UFD；做不了除法就看每个理想是否都是主理想（PID）；再不行才退回用因子链条件与素性条件直接验证 UFD。

\subsection{参考对照}


\begin{itemize}
\item [DF] §8.1（ED 与带余除法）、§8.2（PID）、§8.3（UFD 与两个等价条件）。
\item [DS] §10.2（理想、主理想、$\mathbb Z$ 中理想皆主理想）。该书没有 ED/PID/UFD 这条线。
\item 课程讲义 §7.3：因子链条件与素性条件的定义 7.23、定理 7.5 与定理 7.6（UFD 的两个等价条件）、定理 7.7（PID $\Rightarrow$ UFD）、定义 7.21（欧式整区）与定理 7.4（ED $\Rightarrow$ PID）。
\item 原始材料：本地归档（手写原稿与录音转写），不纳入课程仓库。
\end{itemize}

\section{用理想构造域：商环、域扩张与分裂域}


前三节都在讨论"环里的元素怎么分解"。这一节换一个方向：\textbf{用理想构造新的域}。课堂给了两条路线——第一条是把整环"除以"一个极大理想，第二条是取分式域——然后把得到的域用在多项式上，引出分裂域。

\textbf{本节用到的记号：} $R/M$ 为商环；$F,E,K$ 为域；$F[x]$ 为多项式环；$(f)$ 为 $f$ 生成的主理想；$[E:F]$ 为扩域的\textbf{次数}（下文"域扩张与次数"一节定义）；$\operatorname{Frac}(R)$ 为分式域（在下文"分式域：另一条路线"一节构造）；$\mathbb F_p=\mathbb Z/p\mathbb Z$。

\subsection{极大理想给出域，不可约多项式给出极大理想}


\textbf{定理一。} 设 $R$ 是交换含幺环、$M\trianglelefteq R$ 是极大理想，则 $R/M$ 是域。

\textbf{证明。} 由理想对应定理（[Lec 09 §2](../lec09/sections/02-ideals-and-quotient-rings.md)），$R/M$ 的理想与 $R$ 中包含 $M$ 的理想一一对应。$M$ 极大意味着夹在 $M$ 与 $R$ 之间的理想只有 $M$ 和 $R$ 本身，所以 $R/M$ 只有零理想和自身两个理想。而"交换含幺环没有非平凡理想"等价于"它是域"，故 $R/M$ 是域。$\square$

把这条用在多项式环上，就得到一个非常实用的构造。回忆（第 3 节）：$F$ 是域时 $F[x]$ 是欧几里得整环，因而每个理想都是主理想。

\textbf{定理二。} 设 $F$ 是域、$f\in F[x]$ 非零。则

\[
(f)\ \text{是极大理想}\iff f\ \text{是不可约多项式}.
\]

\textbf{证明。} （$\Rightarrow$）反设 $f=gh$，其中 $g,h$ 都不是常数（即都不是单位）。则 $(f)\subsetneq(g)$：包含显然，严格是因为 $\deg g<\deg f$ 意味着 $f$ 不是 $g$ 的常数倍。这与 $(f)$ 极大矛盾。

（$\Leftarrow$）设 $f$ 不可约且 $(f)\subseteq(g)\subseteq F[x]$。因为 $F[x]$ 是 PID，$(g)$ 由某个 $g$ 生成，且 $g\mid f$。由 $f$ 不可约，$g$ 是单位（此时 $(g)=F[x]$）或 $g\approx f$（此时 $(g)=(f)$）。所以中间没有别的理想，$(f)$ 极大。$\square$

\textbf{推论。} $F[x]/(f)$ 是域，当且仅当 $f$ 不可约。

\textbf{例。} 取 $F=\mathbb R$、$f=x^2+1$。它在 $\mathbb R$ 上不可约（判别式 $<0$），所以

\[
\mathbb R[x]/(x^2+1)\cong\mathbb C .
\]

同构把 $x$ 的像送到 $i$。这正是"复数就是给 $\mathbb R$ 添一个 $x^2=-1$ 的解"的严格说法：在商环里 $x^2+1=0$，即 $x$ 扮演了 $i$ 的角色。

\textbf{例。} 取 $F=\mathbb Q$、$f=x^2-2$（在 $\mathbb Q$ 上不可约，因为 $\sqrt2$ 不是有理数），则 $\mathbb Q[x]/(x^2-2)$ 是域，同构于 $\mathbb Q(\sqrt2)=\{a+b\sqrt2:a,b\in\mathbb Q\}$。

\textbf{顺手记下"添一个元素"的记号。} 设 $R'\supseteq R$ 是更大的环、$u\in R'$。记 $R[u]$ 为 \textbf{$R'$ 中包含 $R$ 与 $u$ 的最小子环}，它由 $u$ 的多项式组成：

\[
R[u]=\{r_0+r_1u+r_2u^2+\cdots+r_ku^k:r_i\in R,\ k\ge0\}.
\]

（有限和。）当 $u$ 是某个不可约多项式的根时，$R[u]$ 的加法结构会"降次"——这正是下面 $F[x]/(f)$ 里每个元素都能用 $1,\bar x,\dots,\bar x^{\,n-1}$ 线性表示的原因。分式域、扩域都建立在这个构造之上。

\textbf{例（手写稿给的二次域同构）。} 对无平方因子的整数 $d$，有环同构

\[
\mathbb Z[\sqrt d]\;\cong\;\mathbb Z[x]/(x^2-d).
\]

映射取"把 $x$ 送到 $\sqrt d$"，即 $\varphi\bigl(\sum a_ix^i\bigr)=\sum a_i(\sqrt d)^i$。它满，因为 $\varphi(a+bx)=a+b\sqrt d$ 已取遍 $\mathbb Z[\sqrt d]$；核由带余除法算出：除以首一多项式 $x^2-d$ 得 $f(x)=q(x)(x^2-d)+r(x)$，$\deg r\le1$，于是 $\varphi(f)=0\iff r(\sqrt d)=0$。若 $r=\alpha+\beta x$ 且 $\alpha+\beta\sqrt d=0$，由 $\sqrt d$ 不是有理数得 $\alpha=\beta=0$，故 $\ker\varphi=(x^2-d)$。用第一同构定理即得。

注意这里 $\mathbb Z[x]/(x^2-d)$ 是\textbf{环}但不一定是域（例如 $x^2-10$ 在 $\mathbb Q$ 上不可约，可 $\mathbb Z[x]/(x^2-10)$ 不是域，因为 $\mathbb Z$ 不是域）。要拿到域，系数环必须是域：$\mathbb Q[x]/(x^2-d)$ 才是域。

\subsection{分式域：另一条路线（回顾）}


如果起点是\textbf{整环}而不是域上的多项式环，可以用分式域把除法"补"进去。$R$ 为整环时，在 $R\times(R\setminus\{0\})$ 上定义

\[
(a,b)\sim(c,d)\iff ad=bc,
\]

把等价类记为 $a/b$，按 $\frac ab+\frac cd=\frac{ad+bc}{bd}$、$\frac ab\cdot\frac cd=\frac{ac}{bd}$ 运算，得到\textbf{分式域} $\operatorname{Frac}(R)$。它满足泛性质：任何把 $R$ 单射进域 $K$ 的同态 $\varphi$，都能唯一地延拓为 $\operatorname{Frac}(R)\to K$ 的同态，取值为

\[
\widetilde{\varphi}(a/b)=\varphi(a)\varphi(b)^{-1}.
\]

用一句话说：\textbf{$\operatorname{Frac}(R)$ 是包含 $R$ 的最小域。} 经典例子是 $\operatorname{Frac}(\mathbb Z)=\mathbb Q$ 与 $\operatorname{Frac}(F[x])=F(x)$（有理函数域）。详细证明见 [Lec 09 §4](../lec09/sections/04-units-fractions-characteristic.md)。

两条路线的分工：\textbf{$F[x]/(f)$ 用来"添一个根"，$\operatorname{Frac}(R)$ 用来"添分母"。} 前者要求 $f$ 不可约，后者要求 $R$ 是整环。

\subsection{域扩张与次数}


若域 $E$ 包含域 $F$（$F\subseteq E$ 且两者的加法乘法一致），就说 $E$ 是 $F$ 的\textbf{扩域}，记作 $E/F$。此时有一个关键观察：

\textbf{$E$ 可以看作 $F$ 上的线性空间。} 向量加法就是 $E$ 的加法；标量乘法取 $F$ 中元素去乘 $E$ 中元素——因为 $F\subseteq E$，这个乘法本来就有，且满足线性空间的所有公理。这个线性空间的维数记作

\[
[E:F]=\dim_F E,
\]

称为扩域的\textbf{次数}。$[E:F]<\infty$ 时说 $E/F$ 是\textbf{有限扩张}。

\textbf{例。} $[\mathbb C:\mathbb R]=2$，因为 $\{1,i\}$ 是 $\mathbb R$-基。$[\mathbb Q(\sqrt2):\mathbb Q]=2$，基为 $\{1,\sqrt2\}$。更一般地，若 $f$ 是 $F$ 上 $n$ 次不可约多项式，则 $E=F[x]/(f)$ 满足 $[E:F]=n$，基为 $\{1,\bar x,\dots,\bar x^{n-1}\}$（$\bar x$ 是 $x$ 在商环中的像）。

\textbf{这解释了 4.1 的构造为什么好用：} 想造一个 $F$ 的 $n$ 次扩域，只要找一个 $F$ 上的 $n$ 次不可约多项式，取 $F[x]/(f)$ 即可。

\textbf{顺带一个计数结论。} $E=F[x]/(f)$ 中每个元素都有唯一的代表元 $\bar a_0+\bar a_1\bar x+\cdots+\bar a_{n-1}\bar x^{\,n-1}$（用带余除法把次数降到 $<n$），所以

\[
|E|=|F|^{\deg f}.
\]

当 $F=\mathbb F_p$、$\deg f=n$ 时就得到 $|E|=p^n$——这正是构造 $p^n$ 元有限域的思路：\textbf{找一个 $\mathbb F_p$ 上的 $n$ 次不可约多项式}。

\subsection{特征与素子域}


域的特征定义为

\[
\operatorname{char}(F)=\min\{n\ge1:\underbrace{1+\cdots+1}_{n}=0\},
\]

若不存在这样的 $n$ 则规定为 $0$。因为域是整环，\textbf{特征不是 $0$ 就一定是素数}（[Lec 09 §4](../lec09/sections/04-units-fractions-characteristic.md)）。

考察所有"若干个 $1$ 相加或相减"得到的元素，它们构成 $F$ 的一个子环：特征为 $p$ 时这个子环同构于 $\mathbb Z/p\mathbb Z=\mathbb F_p$；特征为 $0$ 时同构于 $\mathbb Z$。当 $F$ 是域时，特征 $0$ 的情形还可以再取分式域，于是 $F$ 含有一个同构于 $\mathbb Q$ 的子域（因为 $\mathbb Z$ 的分式域是 $\mathbb Q$）。总结：

\[
\text{特征为 }p:\ F\ \text{含子域}\ \mathbb F_p;\qquad
\text{特征为 }0:\ F\ \text{含子域}\ \mathbb Q .
\]

这个\textbf{最小的子域}（任何子域都必须包含它）称为 $F$ 的\textbf{素子域}；[DF] 把它定义为"由 $1_F$ 生成的子域"，中文教材则按"最小性"把它称为\textbf{素域}（$\mathbb Z_p$ 与 $\mathbb Q$ 统称素域）。两种说法指的是同一个对象。

\textbf{有限域的大小。} 设 $F$ 是有限域，则它的特征不能是 $0$（否则含 $\mathbb Q$，已经无限），所以 $\operatorname{char}(F)=p$ 为素数，且 $F$ 含有素子域 $\mathbb F_p$。于是在 4.3 的意义下 $F$ 是 $\mathbb F_p$ 上的有限维线性空间；设 $[F:\mathbb F_p]=n$，则 $F$ 作为 $\mathbb F_p$-向量空间有 $n$ 个基向量，每个基向量的系数有 $p$ 种选法，因此

\[
|F|=p^{\,n}.
\]

\textbf{结论：有限域的阶必为素数幂。} 反过来的问题——对每个素数幂 $p^n$ 是否都存在 $p^n$ 元域，且同构意义下唯一——正是下一节的内容。

\subsection{分裂域}


设 $f\in F[x]$ 是正次数多项式。若扩域 $E\supseteq F$ 满足

\[
f(x)=c\,(x-\alpha_1)(x-\alpha_2)\cdots(x-\alpha_n),\qquad c\in F^\times,\ \alpha_i\in E,
\]

就说 $f$ 在 $E$ 上\textbf{完全分裂}（split），并把每个 $\alpha_i$ 称为 $f$ 在 $E$ 中的根。使 $f$ 完全分裂的"最小"扩域称为 $f$ 的\textbf{分裂域}（splitting field），即 $E=F(\alpha_1,\dots,\alpha_n)$——由 $F$ 和 $f$ 的全部根一起生成的域。

\textbf{"最小"是定义的一部分，不能漏。} 准确的说法是：$f$ 在 $E$ 上完全分裂，且在\textbf{任何}包含 $F$ 的真子域上都\textbf{不}完全分裂。只要求"$f$ 在 $E$ 上分裂"是不够的（把 $E$ 换成更大的域仍然分裂）。

\textbf{次数上界。} 若 $\deg f=n$，则 $[E:F]\le n!$：每添一个根，扩域次数至多等于当前多项式的次数（4.1 的构造给出 $[F_1:F]=\deg g\le n$），而根的总数不超过 $n$，逐次相乘即得。

\textbf{先记住一个反复用到的等价：根 $\iff$ 一次因子。} 设 $E\supseteq F$ 是扩域、$\alpha\in E$，则

\[
f(\alpha)=0\iff(x-\alpha)\mid f(x)\ \text{在}\ E[x]\ \text{中}.
\]

（$\Leftarrow$ 显然；$\Rightarrow$ 用带余除法除以 $x-\alpha$，余式是常数且等于 $f(\alpha)$。）所以"$f$ 在 $E$ 上完全分裂"等价于"$f$ 在 $E[x]$ 中分解成一次因子之积"，也等价于"$f$ 的\textbf{全部}根都在 $E$ 中"。手写稿把这件事写成"$f$ 可解 $\iff$ $f$ 在 $E$ 中可约（能分解）"。

\textbf{定理（存在性）。} 对任意域 $F$ 与任意正次数多项式 $f\in F[x]$，$f$ 的分裂域存在。

\textbf{证明提纲（对 $\deg f$ 归纳）。} 若 $f$ 的所有不可约因子都是一次，则 $f$ 已在 $F$ 上分裂，$E=F$。否则取 $f$ 的一个次数 $\ge2$ 的不可约因子 $g$，则 $F_1=F[x]/(g)$ 是域（定理二）且 $[F_1:F]=\deg g\ge2$；在 $F_1$ 中 $g$ 有一个根 $\bar x$，所以 $f$ 至少多出一个根。把 $f$ 看作 $F_1[x]$ 中的多项式，它在 $F_1$ 上的不可约因子次数严格变小，于是对 $F_1$ 与 $f$ 用归纳假设，得到 $f$ 在某个 $E\supseteq F_1$ 上完全分裂。$\square$

\textbf{"$\bar x$ 是 $g$ 的根"这句话的验证。} 设 $g(x)=\sum a_ix^i$，$\bar x$ 是 $x$ 在 $F_1=F[x]/(g)$ 中的像。在商环里计算 $g(\bar x)=\sum a_i\bar x^{\,i}=\overline{g(x)}=\bar 0$，因为 $g\in(g)$ 在商环中正是零元。所以 $F_1$ 中确实有了 $g$ 的一个根，从而 $f$ 在这个扩域里多分裂出一个一次因子。

\textbf{唯一性证明的关键一步。} 设 $E,E'$ 都含 $F$ 且 $f$ 在两者中都完全分裂。取 $f$ 的根 $u\in E$，它在 $E'$ 中的对应根为 $u'$。则由 $F$ 与 $u$ 生成的子域 $F(u)$ 到 $F(u')$ 有一个保持 $F$ 不动的同态，把 $u$ 送到 $u'$：因为 $u$ 与 $u'$ 在 $F$ 上有\textbf{同一个极小多项式}（都是 $f$ 的不可约因子）——同态只能把根送到根，而 $F[x]/(p)\to E'$、$\bar x\mapsto u'$ 正是这样的同态。把这个步骤对每个根重复，就得到 $E\to E'$ 的 $F$-同构。$\square$

\textbf{定理（唯一性）。} 分裂域在保持 $F$ 不动的同构意义下唯一：若 $E$ 与 $E'$ 都是 $f$ 在 $F$ 上的分裂域，则存在域同构 $\sigma:E\to E'$ 使 $\sigma|_F=\mathrm{id}_F$。

唯一性要读准：\textbf{"唯一"指的是同构唯一，不是集合相等。} 例如 $f=x^2-2$ 在 $\mathbb Q$ 上，$\mathbb Q(\sqrt2)$ 与 $\mathbb Q(-\sqrt2)$ 作为 $\mathbb Q$ 的子域其实是同一个域；但若取两个"形式不同"的构造，它们之间存在保持 $\mathbb Q$ 不动的同构，把 $\sqrt2$ 送到 $\sqrt2$ 或 $-\sqrt2$。课堂上强调"分类肯定不唯一，但同构意义下唯一"就是这个意思。

\textbf{例。} $f=x^2-2$ 在 $\mathbb Q$ 上的分裂域是 $\mathbb Q(\sqrt2)$，因为两个根 $\pm\sqrt2$ 都在其中生成好了：$(-\sqrt2)=-1\cdot\sqrt2$，而 $-1\in\mathbb Q$。相比之下，$f=x^3-2$ 在 $\mathbb Q$ 上的分裂域比 $\mathbb Q(\sqrt[3]{2})$ 大——后者只含一个实根，另外两个是复根，需要再添 $\omega=e^{2\pi i/3}$。

\textbf{下一节预告。} 有限域的存在唯一性（对每个素数幂 $p^n$ 存在唯一的 $p^n$ 元域 $\mathbb F_{p^n}$）将用分裂域证明：$\mathbb F_{p^n}$ 正是 $x^{p^n}-x$ 在 $\mathbb F_p$ 上的分裂域。

\subsection{参考对照}


\begin{itemize}
\item [DF] §13.1（域扩张与次数）、§13.4（分裂域的存在唯一性）、§14.3（有限域）；§9.3（分式域的一般构造）、§8.2（$F[x]$ 的理想与不可约多项式）。
\item [DS] §10.3（特征、素域、定理 10.3.4 $F[x]/(f)$ 是域 $\iff f$ 不可约、定理 10.3.7 有限域的存在唯一性）。该书的极大/素理想只在习题 10.2\#7，且没有域扩张与分裂域。
\item 课程讲义 §8.1（域扩张与次数 8.1、素子域 8.1、有限域大小为素数幂）、§8.2（分裂域的存在唯一 8.4）、§8.3（有限域 8.5）。
\item 原始材料：../../../课堂手记和录音转写/Lec10/。
\end{itemize}


\section*{一页式总览}
\begin{itemize}
\item 单位群 $R^\times$；$a\mid b$ 是"存在 $c$ 使 $b=ac$"；$a\approx b$ 是"存在单位 $u$ 使 $a=ub$"。
\item 整环中 $a\mid b$ 且 $b\mid a$ $\iff$ $a\approx b$ $\iff$ $(a)=(b)$。所以唯一分解只能要求"在相伴与次序意义下唯一"。
\item 素元：$p\mid ab\Rightarrow p\mid a$ 或 $p\mid b$。不可约元：$p=ab\Rightarrow a$ 或 $b$ 是单位。素元必不可约，反之不然。
\item UFD $\iff$ 因子链条件（无无穷真因子链）$+$ 素性条件（不可约必素）。
\item 域 $\subsetneq$ ED $\subsetneq$ PID $\subsetneq$ UFD $\subsetneq$ 整环；$\mathbb Z$ 与 $F[x]$ 都是 ED，$\mathbb Z[x]$ 是 UFD 但不是 PID。
\item ED：存在范数 $N$ 使带余除法成立；ED $\Rightarrow$ PID 的证明取理想中范数最小的元素。
\item $R$ 交换含幺、$M$ 极大 $\Rightarrow$ $R/M$ 是域；$F[x]/(f)$ 是域 $\iff$ $f$ 不可约。
\item $E/F$ 是扩域 $\Rightarrow$ $E$ 是 $F$-线性空间，$[E:F]$ 是其次数；$|F|=p^n$（有限域）。
\item 分裂域：任何正次数多项式都有分裂域，且在保持 $F$ 不动的同构意义下唯一。
\end{itemize}

\section*{参考资料}
\begin{itemize}
\item \textbf{[DF]} D. S. Dummit, R. M. Foote. \emph{Abstract Algebra}, 3rd ed. §8.1（欧几里得整环）、§8.2（主理想整环）、§8.3（唯一分解整环）、§9.3（分式域）、§13.1--13.2（域扩张与分裂域）、§13.4（有限域）。
\item \textbf{[DS]} 邓小铁、王畅。《离散数学与结构》§10.1（整环）、§10.2（理想与主理想；极大/素理想只在习题 10.2\#7）、§10.3（特征、素域、$F[x]/(f)$ 是域 $\iff f$ 不可约、有限域的存在唯一性）。该书没有 ED/PID/UFD 这条线，也没有域扩张与分裂域。
\item \textbf{[讲义]} 课程讲义 \texttt{25-26-lecture\_notes.pdf} 第 7 章（环）与第 8 章（域）：§7.3 的定理 7.5 与 7.6 给出 UFD 的两个等价条件（这条在 [DF] 中没有以定理形式出现），§8.1--8.3 给出域扩张、分裂域与有限域。该讲义为 2025--2026 秋季版本。
\item 原始材料：\texttt{DiscreteCourse2026/课堂手记和录音转写/Lec10/}（三段转写与手写原稿）。
\end{itemize}
\end{document}